\documentclass{article}
\usepackage[utf8]{inputenc}
\usepackage[T1]{fontenc}
\usepackage{geometry}
\usepackage{graphicx}
\usepackage{float}
\usepackage{hyperref}
\usepackage[english,indonesian]{babel}

\title{Optimasi Pemrosesan Template LaTeX Menggunakan Pendekatan Berbasis AI}
\author{Academic Support Agent}
\date{\today}

\begin{document}
\maketitle

\begin{abstract}

\end{abstract}

\chapter{PENDAHULUAN}


\chapter{TINJAUAN PUSTAKA}


% Original content (commented out):
% \documentclass[12pt]{article}
\usepackage{fullpage}
\usepackage{times}
\usepackage[normalem]{ulem}
\usepackage{fancyhdr,graphicx,amsmath,amssymb, mathtools, scrextend, titlesec, enumitem}
\usepackage[ruled,vlined]{algorithm2e} 
\include{pythonlisting}

\title{Problem 4 Oscillating subsequence}
\author{Luoqi (Rocky) Wu}
\begin{document}
\maketitle

\textbf{Analysis:}

\textbf{Recurrence:}


\pagebreak

\begin{algorithm}
\caption{}
\KwIn{}
\KwOut{}
\hrulefill\\

\nl \tcp*[r]{}
\nl 

\nl \For{i = 2 to k} {
	\nl \For{$j = 0 to 1$}{
    	\nl }
        }
\nl \Return{}
\end{algorithm}

\pagebreak

\textbf{Runtime Analysis:} To bound the running time， \\

\textbf{Proof of Correctness using induction:} \\

Claim: \\

Base Case: \\

Inductive Case: \\

-If \\

-If \\

We have now exhausted the sample space, \\

\end{document}

\end{document}